\chapter{Random Variables and Distribution Functions}\label{ch:random-variables}

\chapterinfo{STAT2003, STAT1301}{\cref{ch:probability-spaces}}

% Sources: STAT2003 handwritten notes (pp. 7-9); STAT2003 course notes, Sections 2.1-2.2;
% STAT1301 course notes (2026), Section 1.8.

Specifying a model for a random experiment via a complete description of $\Omega$ and $\Prob$ is not always convenient or necessary. In practice we are only interested in certain observations (numerical measurements) of the experiment, and we include these in the model via random variables. This chapter defines random variables and their distributions, which are specified by a cumulative distribution function or, in the discrete and continuous cases, by a mass or density function.

% ==========================================================================
\section{Random variables}\label{sec:random-variables}
% ==========================================================================

\begin{definition}[Random variable]\label{def:random-variable}
Let $(\Omega, \mathcal{F}, \Prob)$ be a probability space. A \emph{random variable} is a function $X \colon \Omega \to E$, where $E$ is a measurable space, for example $\R$, a subset of $\R$, $\R^n$, or a space of functions. In this chapter, $E = \R$.
\end{definition}

For a set $A \subseteq \R$, we write
\[
  \set{X \in A} \coloneqq \set{\omega \in \Omega : X(\omega) \in A},
\]
and similarly $\set{X = x}$, $\set{X \le x}$ and $\set{a < X \le b}$. We abbreviate $\Prob(\set{X \in A})$ to $\Prob(X \in A)$. Random variables are usually denoted by capital letters from the end of the alphabet, such as $X, X_1, X_2, \ldots, Y, Z$. They give natural notation for events such as $\set{X = 10}$, $\set{X > 1000}$ and $\set{\max(X, Y) \le Z}$.

\begin{example}[Sum of two dice]\label{ex:sum-two-dice}
Toss two dice one after the other and note their sum. The sample space is $\Omega = \set{(1,1), \ldots, (6,6)}$. The function $X$ defined by $X(i, j) = i + j$ is a random variable, which maps the outcome $(i,j)$ to the sum $i + j$. The set of all outcomes with sum $8$ is $\set{X = 8}$. It contains $5$ outcomes, and all $36$ outcomes are equally likely, so
\[
  \Prob(X = 8) = \frac{5}{36}.
\]
In the same way,
\[
  \Prob(X = x) = \frac{6 - \abs{7 - x}}{36}, \qquad x = 2, \ldots, 12,
\]
which in table form is
\begin{center}
\begin{tabular}{@{}l*{11}{c}@{}}
\toprule
$x$ & 2 & 3 & 4 & 5 & 6 & 7 & 8 & 9 & 10 & 11 & 12 \\
\midrule
$\Prob(X = x)$ & $\frac{1}{36}$ & $\frac{2}{36}$ & $\frac{3}{36}$ & $\frac{4}{36}$ & $\frac{5}{36}$ & $\frac{6}{36}$ & $\frac{5}{36}$ & $\frac{4}{36}$ & $\frac{3}{36}$ & $\frac{2}{36}$ & $\frac{1}{36}$ \\
\bottomrule
\end{tabular}
\end{center}
\Cref{fig:rv-sum-two-dice} depicts $X$ as a function: all outcomes in the encircled set are mapped to 8.
\end{example}

\begin{figure}[H]
\centering
\begin{tikzpicture}[scale=0.55]
  \foreach \i in {1,...,6} \foreach \j in {1,...,6} \fill (\i,\j) circle (2.5pt);
  \foreach \i in {1,...,6} { \node[font=\scriptsize] at (\i,0.3) {\i}; \node[font=\scriptsize] at (0.3,\i) {\i}; }
  \draw[thick, rotate around={-45:(4,4)}, color=thmcolor] (4,4) ellipse (3.3 and 0.55);
  \node at (0.2,6.9) {$\Omega$};
  \draw[thick] (9,3) -- (20,3);
  \foreach \x in {2,...,12} { \draw (8+\x,2.85) -- (8+\x,3.15); \node[font=\scriptsize] at (8+\x,2.4) {\x}; }
  \fill[thmcolor] (16,3) circle (4pt);
  \draw[-{Stealth}, thick, color=thmcolor] (6.4,5.2) to[bend left=25] node[above, font=\small] {$X$} (15.9,3.3);
\end{tikzpicture}
\caption{The random variable $X(i,j) = i + j$ for the sum of two dice. The encircled outcomes form the event $\set{X = 8}$.}
\label{fig:rv-sum-two-dice}
\end{figure}

Although random variables are, mathematically speaking, functions, it is often convenient to view a random variable as an observation of a random experiment that has not yet been carried out: a measurement that becomes available once we carry out the experiment, say tomorrow. The probabilities pertaining to it can be specified today.

\begin{example}[Coin flips]\label{ex:number-of-heads}
Flip a coin $n$ times, where each flip lands heads with probability $p$. Suppose we are not interested in the complete outcome, such as $(0, 1, 0, 1, 1, 0, \ldots)$, but only in the total number of heads (1s). Let $X$ be the function that assigns to each outcome $\omega = (x_1, \ldots, x_n)$ the number $x_1 + \cdots + x_n$. Then $X$ is a random variable, and the event that there are exactly $k$ heads in $n$ flips is $\set{X = k}$. With the coin-tossing model (\cref{sec:coin-tossing-models}),
\[
  \Prob(X = k) = \binom{n}{k} p^k (1 - p)^{n-k}, \qquad k = 0, 1, \ldots, n.
\]
\end{example}

Other examples of random variables are the number of bugs in a computer program and the amount of rain in Brisbane in June.

The set of all values a random variable can take is its \emph{range}. We distinguish between two kinds of random variable:
\begin{itemize}
  \item \emph{discrete} random variables, which can only take isolated (countably many) values. For example, a count can only take non-negative integer values.
  \item \emph{continuous} random variables, which can take any value in an interval or a union of disjoint intervals, such as the positive real line, and have $\Prob(X = c) = 0$ for every possible value $c$. For example, rainfall, lifetimes and lengths are, at least in principle, continuous.
\end{itemize}

% ==========================================================================
\section{Probability distributions}\label{sec:distribution}
% ==========================================================================

\begin{definition}[Probability distribution]\label{def:distribution}
The \emph{probability distribution} of a random variable $X$ is a description of the probabilities
\[
  \Prob(X \in A) = \Prob(\set{\omega \in \Omega : X(\omega) \in A}), \qquad A \subseteq \R.
\]
If we can specify all probabilities involving $X$, we have specified its distribution.
\end{definition}

The description may take the form of a cumulative distribution function, a probability mass function, a probability density function, a transform (\cref{ch:generating-functions}), and so on.

% ==========================================================================
\section{The cumulative distribution function}\label{sec:cdf}
% ==========================================================================

One way to specify the distribution of $X$ is to give the probabilities of all events of the form $\set{X \le x}$, $x \in \R$.

\begin{definition}[Cumulative distribution function]\label{def:cdf}
The \emph{cumulative distribution function} (CDF) of a random variable $X$ is the function $F \colon \R \to [0, 1]$ defined by
\[
  F(x) \coloneqq \Prob(X \le x), \qquad x \in \R.
\]
We write $F_X$ when the random variable needs naming.
\end{definition}

\Cref{fig:generic-cdf} shows the graph of a CDF.

\begin{figure}[H]
\centering
\begin{tikzpicture}
\begin{axis}[bookplot, width=0.6\textwidth, height=4.2cm, xmin=-3, xmax=4, ymin=0, ymax=1.1, xtick=\empty, ytick={0,1}, xlabel={$x$}, ylabel={$F(x)$}]
  \addplot[no marks, color=thmcolor, domain=-3:0, samples=40] {0.3/(1+exp(-2*(x+1)))};
  \addplot[no marks, color=thmcolor, domain=0:1] {0.6};
  \addplot[no marks, color=thmcolor, domain=1:4, samples=40] {1 - 0.4*exp(-1.5*(x-1))};
  \addplot[only marks, mark=*, mark size=1.5pt, color=thmcolor] coordinates {(0,0.6)};
  \addplot[only marks, mark=*, mark size=1.5pt, mark options={fill=white}, color=thmcolor] coordinates {(0,0.2647)};
  \addplot[dashed, thin, domain=-3:4] {1};
\end{axis}
\end{tikzpicture}
\caption{A cumulative distribution function: non-decreasing, between 0 and 1, and right-continuous (filled dot) at a jump.}
\label{fig:generic-cdf}
\end{figure}

\begin{theorem}[Properties of a CDF]\label{thm:cdf-properties}
Let $F$ be the CDF of a random variable $X$. Then
\begin{enumerate}[label=(\roman*)]
  \item $F$ is non-decreasing: $x \le y \implies F(x) \le F(y)$
  \item $0 \le F(x) \le 1$ for all $x$
  \item $\lim_{x \to -\infty} F(x) = 0$ and $\lim_{x \to \infty} F(x) = 1$
  \item $F$ is right-continuous: $\lim_{h \downarrow 0} F(x + h) = F(x)$
\end{enumerate}
\end{theorem}

\begin{proof}\leavevmode
\begin{enumerate}[leftmargin=*]
  \item[(i)] If $x \le y$, then $A = \set{X \le x} \subseteq B = \set{X \le y}$, so $\Prob(A) \le \Prob(B)$ (\cref{thm:probability-properties}).
  \item[(ii)] $0 \le \Prob(A) \le 1$ for every event $A$, in particular for $A = \set{X \le x}$.
\end{enumerate}
Properties (iii) and (iv) are not proved in this course.
\end{proof}

CDFs are useful because they apply to both discrete and continuous distributions. Any function $F$ with the above properties can be used to specify the distribution of a random variable.

\begin{proposition}[Probabilities from the CDF]\label{prop:cdf-interval-probabilities}
Let $X$ have CDF $F$, and let $a < b$. Then
\begin{enumerate}[label=(\alph*)]
  \item $\Prob(a < X \le b) = F(b) - F(a)$
  \item $\Prob(X = a) = F(a) - \lim_{h \downarrow 0} F(a - h)$
  \item $\Prob(a \le X \le b) = F(b) - \lim_{h \downarrow 0} F(a - h)$
\end{enumerate}
\end{proposition}

\begin{proof}\leavevmode
\begin{enumerate}[leftmargin=*]
  \item[(a)] $\set{X \le b} = \set{X \le a} \cup \set{a < X \le b}$, and these two events are disjoint. By the sum rule, $F(b) = F(a) + \Prob(a < X \le b)$.
  \item[(b)] Not proved in this course.
  \item[(c)] By (a) and (b),
  \[
    \Prob(a \le X \le b) = \Prob(a < X \le b) + \Prob(X = a) = F(b) - \lim_{h \downarrow 0} F(a - h). \qedhere
  \]
\end{enumerate}
\end{proof}

% ==========================================================================
\section{Discrete distributions}\label{sec:discrete-distributions}
% ==========================================================================

\begin{definition}[Discrete distribution, probability mass function]\label{def:discrete-distribution}
A random variable $X$ has a \emph{discrete distribution} if $X$ is a discrete random variable. In particular, for some finite or countable set of values $x_1, x_2, \ldots$,
\[
  \Prob(X = x_i) > 0 \quad (i = 1, 2, \ldots) \qquad \text{and} \qquad \sum_i \Prob(X = x_i) = 1.
\]
The \emph{probability mass function} (pmf) $f$ of $X$ is $f(x) = \Prob(X = x)$. We write $f_X$ to stress that the pmf belongs to $X$.
\end{definition}

The easiest way to specify the distribution of a discrete random variable is to give its pmf. If we know $f(x)$ for all $x$, the sum rule gives every probability involving $X$:
\begin{equation}\label{eq:pmf-sum-rule}
  \Prob(X \in B) = \sum_{x \in B} f(x)
\end{equation}
for any subset $B$ of the range of $X$ (say an interval $(a, b)$), as illustrated in \cref{fig:pmf-sum-rule}.

\begin{figure}[H]
\centering
\begin{tikzpicture}
\begin{axis}[pmfplot, width=0.6\textwidth, height=4.2cm, xmin=0, xmax=11, ymin=0, ymax=0.25, xtick=\empty, ytick=\empty, xlabel={$x$}, ylabel={$f(x)$}, clip=false, xlabel style={at={(1,0)}, anchor=north east, yshift=-2pt}]
  \fill[defcolor, opacity=0.12] (axis cs:3.5,0) rectangle (axis cs:7.5,0.24);
  \addplot[ycomb, mark=*, mark size=1.5pt, color=remcolor] coordinates {(1,0.05) (2,0.12) (3,0.18) (8,0.08) (9,0.06) (10,0.03)};
  \addplot[ycomb, mark=*, mark size=1.8pt, color=defcolor, very thick] coordinates {(4,0.2) (5,0.15) (6,0.1) (7,0.03)};
  \draw[decorate, decoration={brace, mirror, amplitude=5pt}, thick, defcolor] (axis cs:3.6,-0.008) -- node[below=5pt] {$B$} (axis cs:7.4,-0.008);
\end{axis}
\end{tikzpicture}
\caption{A probability mass function. $\Prob(X \in B)$ is the total of the highlighted masses $f(x)$, $x \in B$.}
\label{fig:pmf-sum-rule}
\end{figure}

\begin{remark}\label{rem:pmf-from-cdf}
By \cref{prop:cdf-interval-probabilities}(b), the pmf can be recovered from the CDF:
\[
  f(x) = \Prob(X = x) = F(x) - \lim_{u \uparrow x} F(u).
\]
If $X$ takes only integer values, then $\lim_{u \uparrow x} F(u) = F(x - 1)$ for integer $x$, so $f(x) = F(x) - F(x - 1)$. This simpler formula does not hold for a discrete random variable with non-integer values.
\end{remark}

\begin{example}[Tossing a die]\label{ex:tossing-a-die}
Toss a die and let $X$ be its face value. $X$ is discrete with range $\set{1, 2, \ldots, 6}$. If the die is fair, the pmf is
\begin{center}
\begin{tabular}{@{}l*{6}{c}c@{}}
\toprule
$x$ & 1 & 2 & 3 & 4 & 5 & 6 & $\sum$ \\
\midrule
$f(x)$ & $\frac16$ & $\frac16$ & $\frac16$ & $\frac16$ & $\frac16$ & $\frac16$ & 1 \\
\bottomrule
\end{tabular}
\end{center}
\end{example}

\begin{example}[Largest value of two dice]\label{ex:max-two-dice}
Toss two dice and let $X$ be the largest face value showing. The pmf of $X$ is
\begin{center}
\begin{tabular}{@{}l*{6}{c}c@{}}
\toprule
$x$ & 1 & 2 & 3 & 4 & 5 & 6 & $\sum$ \\
\midrule
$f(x)$ & $\frac{1}{36}$ & $\frac{3}{36}$ & $\frac{5}{36}$ & $\frac{7}{36}$ & $\frac{9}{36}$ & $\frac{11}{36}$ & 1 \\
\bottomrule
\end{tabular}
\end{center}
Find the probability that the maximum is at least $3$.
\end{example}

\begin{solution}
From the table and \eqref{eq:pmf-sum-rule}, $\Prob(X \ge 3) = \sum_{x=3}^{6} f(x) = \frac{32}{36} = \frac89$.
\end{solution}

% ==========================================================================
\section{Continuous distributions}\label{sec:continuous-distributions}
% ==========================================================================

\begin{definition}[Continuous distribution, probability density function]\label{def:continuous-distribution}
A random variable $X$ has a \emph{continuous distribution} if $X$ is a continuous random variable for which there exists a function $f \colon \R \to \R$ with $f(x) \ge 0$ for all $x$ and $\int_{-\infty}^{\infty} f(x) \dd x = 1$, such that for all $a \le b$,
\[
  \Prob(a < X \le b) = F(b) - F(a) = \int_a^b f(u) \dd u.
\]
The function $f$ is the \emph{probability density function} (pdf) of $X$.
\end{definition}

We write $X \sim f$ or $X \sim F$ to indicate that $X$ is distributed according to the pdf $f$ or the CDF $F$. The same notation is used for a pmf. Once we know the pdf, we can calculate the probability that $X$ lies in a set $B$ by integration, in analogy with \eqref{eq:pmf-sum-rule}:
\[
  \Prob(X \in B) = \int_B f(x) \dd x,
\]
as illustrated in \cref{fig:pdf-integral-rule}.

\begin{figure}[H]
\centering
\begin{tikzpicture}
\begin{axis}[bookplot, width=0.6\textwidth, height=4.2cm, xmin=-1, xmax=6, ymin=0, ymax=0.45, xtick={1.5,3.2}, xticklabels={$a$, $b$}, ytick=\empty, xlabel={$x$}, ylabel={$f(x)$}]
  \addplot[name path=f, no marks, color=thmcolor, domain=-1:6, samples=80] {0.4*exp(-0.5*(x-2)^2/1.2)};
  \path[name path=axis] (axis cs:-1,0) -- (axis cs:6,0);
  \addplot[thmcolor!25] fill between[of=f and axis, soft clip={domain=1.5:3.2}];
\end{axis}
\end{tikzpicture}
\caption{A probability density function. $\Prob(a < X \le b)$ is the shaded area under $f$ between $a$ and $b$.}
\label{fig:pdf-integral-rule}
\end{figure}

The CDF $F$ is a primitive (antiderivative) of the pdf $f$. In particular,
\[
  F(x) = \Prob(X \le x) = \int_{-\infty}^x f(u) \dd u,
\]
and $f$ is the derivative of $F$:
\[
  f(x) = \frac{\dd}{\dd x} F(x) = F'(x).
\]
We can interpret $f(x)$ as the ``density'' that $X = x$. More precisely,
\[
  \Prob(x \le X \le x + h) = \int_x^{x + h} f(u) \dd u \approx h f(x).
\]

\begin{remark}\label{rem:density-not-probability}
$f(x)$ is not a probability: it is a probability density. In particular, if $X$ is continuous, then
\[
  \Prob(X = x) = 0 \qquad \text{for all } x.
\]
This also justifies using $\Prob(x \le X \le x + h)$ above instead of $\Prob(x < X \le x + h)$: for a continuous random variable, including or excluding an endpoint does not change the probability. The same letter $f$ is used for the pmf in the discrete case and the pdf in the continuous case, but it is crucial to understand the difference between the two.
\end{remark}

\begin{example}[Drawing uniformly from an interval]\label{ex:uniform-0-2}
Draw a random number $X$ from the interval $[0, 2]$, where each number is equally possible. Find the pdf $f$ and the CDF $F$ of $X$.
\end{example}

\begin{solution}
Drawing a number ``uniformly'' from $[0, 2]$ means that $\Prob(X \le x) = x/2$ for all $x \in [0, 2]$. So the CDF is
\[
  F(x) = \begin{cases} 0, & x < 0, \\ x/2, & 0 \le x \le 2, \\ 1, & x > 2. \end{cases}
\]
Differentiating $F$ gives
\[
  f(x) = \begin{cases} 1/2, & 0 \le x \le 2, \\ 0, & \text{otherwise}. \end{cases}
\]
The density is constant on $[0, 2]$ and zero elsewhere, which reflects that each point in $[0, 2]$ is equally likely. Moreover, every subinterval of $[0, 2]$ of width $a$, with $0 < a < 2$, is equally likely.
\end{solution}

\begin{definition}[Quantiles and percentiles]\label{def:quantiles}
Let $X$ be a continuous random variable with CDF $F$ and pdf $f$, and let $0 < p < 1$. The \emph{$p$th quantile}, or \emph{$100p$th percentile}, of the distribution of $X$ is the number $q(p)$ with
\[
  p = F\bigl(q(p)\bigr) = \int_{-\infty}^{q(p)} f(u) \dd u.
\]
If $F$ can be inverted, this can be written $q(p) = F^{-1}(p)$, which is called the quantile function.
\end{definition}

Describing an experiment through a random variable and its pmf, pdf or CDF is usually much easier than specifying the probability space $(\Omega, \mathcal{F}, \Prob)$ explicitly. In the examples above, no probability space was needed.
